% ============================================================
%  Chapter 6 — The Implemented System
% ============================================================
\chapter{The Implemented System}
\label{ch:implementacio}

This chapter describes the software built around the KCMamba model:
the trading engine, dashboard, backtesting infrastructure, and Binance adapter.
The system follows a hexagonal architecture (\emph{Ports \& Adapters}); the
architectural principles were introduced in Sec.~\ref{sec:system_arch};
here the focus is on the concrete implementation of each adapter.

% ─────────────────────────────────────────────────────────────────────────────
\section{Technology Stack}
\label{sec:impl_tech}
% ─────────────────────────────────────────────────────────────────────────────

The entire system is written in \textbf{Python 3.11}.

\begin{table}[H]
  \centering
  \small
  \begin{tabular}{|l|l|p{7cm}|}
    \hline
    \textbf{Layer} & \textbf{Technology} & \textbf{Purpose} \\
    \hline\hline
    Model          & PyTorch~2.x          & KCMamba forward pass, gradient computation \\
    I/O            & asyncio / aiohttp    & Event-driven, non-blocking I/O \\
    Exchange       & python-binance       & Live market data, order execution \\
    Web server     & aiohttp (HTTP)       & Static HTML/JS serving, REST JSON endpoints \\
    Frontend       & Vanilla JS + Chart.js & Interactive charts, no build step required \\
    Testing        & pytest + unittest.mock & Unit and integration tests (backtester) \\
    Configuration  & pydantic v2          & Strongly-typed settings validation \\
    Logging        & Python logging       & Rotating file + console handler \\
    \hline
  \end{tabular}
  \caption{Main technology components of the system and their roles.}
  \label{tab:tech_stack}
\end{table}

The frontend deliberately avoids React or Vue: every page is a standalone
static HTML file served directly by \texttt{aiohttp}. No \texttt{npm} build step
is required, yet Chart.js provides fully interactive visualisations.

% ─────────────────────────────────────────────────────────────────────────────
\section{Component Design}
\label{sec:component_plan}
% ─────────────────────────────────────────────────────────────────────────────

The system runs in a single Python process. Component interactions:

\begin{description}
  \item[\textbf{Entry point} (\texttt{main.py})]
        Starts the \texttt{asyncio} event loop and initialises
        \texttt{TradingEngine}, \texttt{MLService}, and the adapter.

  \item[\textbf{TradingEngine} (\texttt{core/application/engine.py})]
        The centre of business logic. Receives \texttt{PriceEvent}s via
        \texttt{InMemoryEventBus}, runs pipeline stages in order
        (Prediction $\to$ Signal $\to$ Risk $\to$ Execution), and issues orders
        through \texttt{IBrokerGatewayPort}.

  \item[\textbf{aiohttp web server} (\texttt{adapters/web/dashboard.py})]
        Runs in the same \texttt{asyncio} event loop as the trading loop.
        Pushes a 1~Hz WebSocket stream to browsers via an \texttt{asyncio.Task};
        also exposes REST endpoints (\texttt{/api/signals}, \texttt{/api/equity},
        etc.) for state queries.

  \item[\textbf{Binance WebSocket / REST} (external)]
        \texttt{BinanceFeed} receives OHLCV data over an async WebSocket;
        \texttt{BinanceBroker} sends orders via REST.
        Automatic reconnection on network interruptions.

  \item[\textbf{Filesystem}]
        Model weights (\texttt{*.pt}), log files, CSV exports, and the
        \texttt{.env} configuration file. No database is used;
        all runtime state lives in in-memory data structures.
\end{description}

\begin{table}[H]
  \centering
  \small
  \begin{tabular}{|l|l|l|l|}
    \hline
    \textbf{Component} & \textbf{Protocol} & \textbf{Direction} & \textbf{Role} \\
    \hline\hline
    TradingEngine $\leftrightarrow$ EventBus   & in-process call & bidirectional & event pub/sub \\
    TradingEngine $\to$ BinanceBroker          & REST HTTPS      & outgoing      & order dispatch \\
    BinanceFeed $\leftarrow$ Binance           & WebSocket WSS   & incoming      & price stream \\
    DashboardServer $\leftrightarrow$ browser  & HTTP + WS       & bidirectional & visualisation \\
    MLService $\leftrightarrow$ filesystem     & file I/O        & bidirectional & weight load/save \\
    \hline
  \end{tabular}
  \caption{Runtime component connectivity.}
  \label{tab:component_plan}
\end{table}

% ─────────────────────────────────────────────────────────────────────────────
\section{The Web Dashboard}
\label{sec:dashboard}
% ─────────────────────────────────────────────────────────────────────────────

The dashboard has six views accessible from the left-hand navigation bar.
The backend pushes updates via a 1~Hz WebSocket stream.

% · · · · · · · · · · · · · · · · · · · · · · · · · · · · · · · · · · · · ·
\subsection{Home: Overview}
\label{subsec:dash_overview}
% · · · · · · · · · · · · · · · · · · · · · · · · · · · · · · · · · · · · ·

The home page (\texttt{index.html}) consolidates the most important information:
equity curve, open positions, recent trading signals, and the model's current
Kalman gain value.

\begin{figure}[H]
  \centering
  \fbox{\begin{minipage}[c][6cm][c]{0.92\textwidth}
    \centering
    \textit{[Screenshot: Home --- Overview]}\\[4pt]
    \small Placeholder for \texttt{dashboard/index.html} screenshot\\
    showing the equity curve, open positions, and Kalman gain.
  \end{minipage}}
  \caption{Dashboard home page. The top panel shows the equity curve; below it
           are open positions and recent signals.}
  \label{fig:dash_overview}
\end{figure}

% · · · · · · · · · · · · · · · · · · · · · · · · · · · · · · · · · · · · ·
\subsection{Trading Signals}
\label{subsec:dash_signals}
% · · · · · · · · · · · · · · · · · · · · · · · · · · · · · · · · · · · · ·

The signal log (\texttt{signals.html}) displays all generated
\textsc{Buy}/\textsc{Sell}/\textsc{Hold} signals with timestamp, symbol,
signal strength, and the model's current forecast value. The view is
filterable by symbol and signal type.

\begin{figure}[H]
  \centering
  \fbox{\begin{minipage}[c][6cm][c]{0.92\textwidth}
    \centering
    \textit{[Screenshot: Trading Signals --- \texttt{signals.html}]}\\[4pt]
    \small Signal log table in chronological order, with filter and export.
  \end{minipage}}
  \caption{Trading signals view. Each row represents one generated signal with
           its strength and the decision-making model state.}
  \label{fig:dash_signals}
\end{figure}

% · · · · · · · · · · · · · · · · · · · · · · · · · · · · · · · · · · · · ·
\subsection{Model Diagnostics}
\label{subsec:dash_model}
% · · · · · · · · · · · · · · · · · · · · · · · · · · · · · · · · · · · · ·

The model diagnostics page (\texttt{model.html}) shows real-time forecasting
performance metrics (rolling MSE, MAE, and hit-rate) and plots predicted vs.\
actual values.

\begin{figure}[H]
  \centering
  \fbox{\begin{minipage}[c][6cm][c]{0.92\textwidth}
    \centering
    \textit{[Screenshot: Model Diagnostics --- \texttt{model.html}]}\\[4pt]
    \small Predicted vs.\ actual values, rolling MSE, confidence interval.
  \end{minipage}}
  \caption{Model diagnostics view. The Chart.js panel updates in real time
           with forecasts and residuals.}
  \label{fig:dash_model}
\end{figure}

% · · · · · · · · · · · · · · · · · · · · · · · · · · · · · · · · · · · · ·
\subsection{Kalman Filter Diagnostics}
\label{subsec:dash_kalman}
% · · · · · · · · · · · · · · · · · · · · · · · · · · · · · · · · · · · · ·

The Kalman diagnostics view (\texttt{kla\_kalman.html}) shows the model's
internal state: the adaptive $K$-gain, forgetting factor $A$, and noise
estimate $R$ over time. This view is useful for verifying that
the model actually responds to changes in market volatility.

\begin{figure}[H]
  \centering
  \fbox{\begin{minipage}[c][6cm][c]{0.92\textwidth}
    \centering
    \textit{[Screenshot: Kalman Filter Diagnostics --- \texttt{kla\_kalman.html}]}\\[4pt]
    \small $K$-gain, $A$-factor, and $R$ noise estimate over time during
    moderate and high-volatility market periods.
  \end{minipage}}
  \caption{Kalman filter diagnostics. Decreasing $K$-gain at high volatility
           confirms that the adaptive filtering mechanism is working correctly.}
  \label{fig:dash_kalman}
\end{figure}

% · · · · · · · · · · · · · · · · · · · · · · · · · · · · · · · · · · · · ·
\subsection{Offline Training}
\label{subsec:dash_training}
% · · · · · · · · · · · · · · · · · · · · · · · · · · · · · · · · · · · · ·

The training view (\texttt{training.html}) monitors offline training progress:
epoch loss curve, validation MSE, estimated completion time, and the
current configuration parameters.

\begin{figure}[H]
  \centering
  \fbox{\begin{minipage}[c][6cm][c]{0.92\textwidth}
    \centering
    \textit{[Screenshot: Offline Training --- \texttt{training.html}]}\\[4pt]
    \small Epoch loss curve, validation MSE with live updates, progress bar,
    and configuration summary.
  \end{minipage}}
  \caption{Offline training dashboard. The full training run can be monitored,
           including early-stopping status.}
  \label{fig:dash_training}
\end{figure}

% · · · · · · · · · · · · · · · · · · · · · · · · · · · · · · · · · · · · ·
\subsection{Benchmark Comparison}
\label{subsec:dash_benchmark}
% · · · · · · · · · · · · · · · · · · · · · · · · · · · · · · · · · · · · ·

The benchmark view (\texttt{ltsf\_benchmark.html}) displays LSTM, Fair Mamba,
and KCMamba results as interactive heat maps and line charts, broken down by
noise level and stride configuration.

\begin{figure}[H]
  \centering
  \fbox{\begin{minipage}[c][6cm][c]{0.92\textwidth}
    \centering
    \textit{[Screenshot: LSTM/Mamba comparison --- \texttt{ltsf\_benchmark.html}]}\\[4pt]
    \small Interactive heat map and line charts for MSE results across noise
    levels and datasets.
  \end{minipage}}
  \caption{Benchmark comparison view. Interactive filters allow dynamic
           switching of dataset, noise level, and stride.}
  \label{fig:dash_benchmark}
\end{figure}

% ─────────────────────────────────────────────────────────────────────────────
\section{Backtesting Framework}
\label{sec:backtest}
% ─────────────────────────────────────────────────────────────────────────────

The backtester (\texttt{adapters/testing/backtest.py}) runs all trading logic
through the same \texttt{IBrokerGatewayPort} and \texttt{IMarketFeedPort}
interfaces as the live system, using a simulated event bus
(\texttt{InMemoryEventBus}) and a deterministic time source
(\texttt{FixedTimeSource}) to ensure reproducibility.

A run loads and normalises historical OHLCV data, trains the model offline,
then steps through the validation period minute by minute,
applying 0.1\% taker and 0.1\% maker fees. Results are logged as an equity curve, Sharpe
ratio, maximum drawdown, and win rate.

\begin{figure}[H]
  \centering
  \fbox{\begin{minipage}[c][5.5cm][c]{0.92\textwidth}
    \centering
    \textit{[Screenshot: Backtest results --- equity curve and metrics]}\\[4pt]
    \small Equity curve, Sharpe ratio, maximum drawdown, and monthly
    return heat map for the backtest period.
  \end{minipage}}
  \caption{Backtest results summary. The equity curve shows the strategy's
           cumulative performance over the historical period.}
  \label{fig:backtest_results}
\end{figure}

% ─────────────────────────────────────────────────────────────────────────────
\section{Risk Management}
\label{sec:risk_management}
% ─────────────────────────────────────────────────────────────────────────────

The risk layer (\texttt{core/domain/risk.py}) enforces four rules.
Position-size cap: no single entry may exceed 10\% of total equity.
Asymmetric stop-loss: $-2\%$ from entry on long positions, $-1.5\%$ on short.
If more than five symbols are open simultaneously, new signals are blocked;
and if the daily realised loss exceeds 3\% of equity, the engine closes all
open positions.

Additionally, an order is only dispatched if the KCMamba $K$-gain exceeds
the threshold $K_\text{min} = 0{.}15$, indicating sufficient model confidence
given the current data quality.

\begin{figure}[H]
  \centering
  \fbox{\begin{minipage}[c][5cm][c]{0.92\textwidth}
    \centering
    \textit{[Screenshot: Risk log and position-sizing panel]}\\[4pt]
    \small Risk module log: attempted and rejected signals, stop-loss event
    sequence, daily loss progress bar.
  \end{minipage}}
  \caption{Risk management overview. Red rows indicate rejected signals;
           green rows indicate executed ones.}
  \label{fig:risk_panel}
\end{figure}

% ─────────────────────────────────────────────────────────────────────────────
\section{Live Trading with Binance}
\label{sec:live_trading}
% ─────────────────────────────────────────────────────────────────────────────

The live trading adapter (\texttt{adapters/infrastructure/binance\_broker.py})
uses the Binance Futures REST and WebSocket endpoints. Price updates reach the
\texttt{TradingEngine} via \texttt{IEventBusPort}; the business logic has no
knowledge of the Binance-specific data format; the adapter is replaceable.

\begin{figure}[H]
  \centering
  \fbox{\begin{minipage}[c][5cm][c]{0.92\textwidth}
    \centering
    \textit{[Screenshot: Live trading --- Binance positions and orders]}\\[4pt]
    \small Open Binance Futures positions, pending orders, and equity updates
    in real time.
  \end{minipage}}
  \caption{Live Binance trading view. The panel is synchronised with the
           Binance account; balance and positions update every second.}
  \label{fig:live_trading}
\end{figure}

The system can also run in \emph{paper trading} mode: the Binance adapter is
replaced with \texttt{adapters/testing/broker.py}'s simulated executor,
allowing the algorithm to be tested on live market data without risking real
capital.

% ─────────────────────────────────────────────────────────────────────────────
\section{Running the System}
\label{sec:running}
% ─────────────────────────────────────────────────────────────────────────────

The entire system is started through a single entry point, \texttt{main.py}:

\begin{verbatim}
# Backtest mode (no network connection required)
python main.py --mode backtest --symbol BTCUSDT --from 2024-01-01 --to 2024-12-31

# Paper trading mode (Binance feed, simulated execution)
python main.py --mode paper --symbol BTCUSDT ETHUSDT

# Live trading mode (Binance API key required)
python main.py --mode live --symbol BTCUSDT
\end{verbatim}

The dashboard starts automatically and is accessible at
\texttt{http://localhost:8080}. API keys are loaded from a \texttt{.env} file
and never appear in source code.

\begin{figure}[H]
  \centering
  \fbox{\begin{minipage}[c][4.5cm][c]{0.92\textwidth}
    \centering
    \textit{[Screenshot: Terminal --- system startup and log lines]}\\[4pt]
    \small Console output during startup: model loading, WebSocket connection
    confirmation, event loop start.
  \end{minipage}}
  \caption{System startup log. Each adapter initialises in sequence; on error
           the system prints a detailed message and stack trace.}
  \label{fig:startup_log}
\end{figure}

% ─────────────────────────────────────────────────────────────────────────────
\section{User Documentation}
\label{sec:user_docs}
% ─────────────────────────────────────────────────────────────────────────────

\subsection{Purpose and Audience}

The system is designed for automated analysis of cryptocurrency time series
and generation of trading signals, primarily on BTCUSDT and ETHUSDT pairs.
Its target audience is developers and researchers who want to backtest or
live-run ML-based trading strategies.

\subsection{System Requirements}

\begin{table}[H]
  \centering
  \small
  \begin{tabular}{|l|l|l|}
    \hline
    \textbf{Component} & \textbf{Minimum} & \textbf{Recommended} \\
    \hline\hline
    Operating system & Linux/macOS/Windows 10+ & Ubuntu 22.04 LTS \\
    Python           & 3.10                   & 3.11 \\
    RAM              & 4~GB                   & 8~GB \\
    Disk             & 2~GB                   & 5~GB \\
    GPU (optional)   & ---                    & NVIDIA $\ge$6~GB VRAM \\
    Network          & --- (backtest mode)    & Stable internet (live) \\
    \hline
  \end{tabular}
  \caption{Minimum and recommended hardware/software requirements.}
  \label{tab:sysreq}
\end{table}

\subsection{Installation Guide}

\begin{enumerate}
  \item Clone and create a virtual environment:
\begin{verbatim}
git clone <repo_url> && cd diplomamunkakod
python -m venv .venv && source .venv/bin/activate
\end{verbatim}
  \item Install dependencies (ML packages are optional):
\begin{verbatim}
pip install -r requirements.txt
pip install -r requirements-ml.txt   # GPU: CUDA-compatible PyTorch
\end{verbatim}
  \item Configure API keys (only required for live/paper mode):
\begin{verbatim}
cp .env.example .env
# Edit .env: BINANCE_API_KEY, BINANCE_SECRET
\end{verbatim}
\end{enumerate}

\subsection{Usage Guide}

The application is controlled via \texttt{main.py}
(see Sec.~\ref{sec:running}). The dashboard is available at
\texttt{http://localhost:8080}; backtest results can be exported as CSV and
logs are stored in \texttt{run\_log.txt}.

\subsection{Troubleshooting}

\begin{description}
  \item[\texttt{ModuleNotFoundError: mamba\_ssm}] Install the CUDA-compatible
       packages from \texttt{requirements-ml.txt}, or run on CPU:
       \texttt{python main.py --device cpu}.
  \item[Binance API 403/401] Check the \texttt{.env} keys and IP whitelist
       settings on the Binance trading interface.
  \item[WebSocket connection error] A firewall may block Binance's
       \texttt{wss://} endpoint (port $9443$); set up a proxy if necessary.
\end{description}

% ─────────────────────────────────────────────────────────────────────────────
\section{Testing}
\label{sec:testing}
% ─────────────────────────────────────────────────────────────────────────────

\subsection{Testing Strategy}

The test suite uses \textbf{pytest}. Core logic is covered by unit tests;
the Binance and aiohttp adapters are replaced with \texttt{unittest.mock}, so
all tests run deterministically without a network connection.

\begin{table}[H]
  \centering
  \small
  \begin{tabular}{|l|l|p{7cm}|}
    \hline
    \textbf{Test file} & \textbf{Type} & \textbf{Coverage} \\
    \hline\hline
    \texttt{test\_strategy.py}           & Unit        & ThresholdStrategy signal logic, threshold/config/scale checks \\
    \texttt{test\_risk\_policy.py}        & Unit        & RiskPolicy blocking paths, MaxQty, MaxNotional, Confidence, SlippageGuard \\
    \texttt{test\_execution.py}          & Integration & PaperBroker order execution, fill simulation, Stop/Limit orders \\
    \texttt{test\_equity.py}             & Unit        & SimulationWallet equity calculation, PnL, mark-price updates \\
    \texttt{test\_fee\_fix.py}           & Regression  & Commission calculation accuracy across a trading cycle \\
    \texttt{test\_fill\_fix.py}          & Regression  & Fill-event time/price logic, partial fill handling \\
    \texttt{test\_signal\_diagnostics.py} & Unit       & Signal diagnostic metrics (edge score, hit-rate calculation) \\
    \texttt{test\_binance\_feed.py}      & Mock-based  & BinanceFeed WebSocket event processing, reconnect logic \\
    \hline
  \end{tabular}
  \caption{Summary of test files.}
  \label{tab:tests}
\end{table}

\subsection{Running the Tests}

\begin{verbatim}
# Run all tests
pytest tests/ -v

# Unit tests only, skipping slow integration tests
pytest tests/ -v -m "not slow"

# Coverage report
pytest tests/ --cov=core --cov-report=term-missing
\end{verbatim}

\subsection{Selected Test Cases}

\paragraph{Risk rules} \texttt{test\_risk\_policy.py} verifies that all
blocking rules (MaxQtyRule, MaxNotionalRule, ConfidenceRule, SlippageGuardRule)
correctly reject invalid orders, and that when multiple rules are applied
together the strictest one prevails. Tests cover approximately 100 random input
combinations.

\paragraph{Strategy signal logic} \texttt{test\_strategy.py} verifies that
\texttt{ThresholdStrategy} emits a \textsc{Buy}/\textsc{Sell} signal exactly
when the predicted probability exceeds \texttt{threshold}, and that
\texttt{min\_edge}, \texttt{min\_confidence}, and \texttt{max\_sigma} correctly
filter weak signals.

\paragraph{Equity and fill logic} All tests in \texttt{test\_equity.py} and
\texttt{test\_fill\_fix.py} pass, confirming that PnL calculation and partial
fill handling work correctly in multi-symbol portfolios.

\paragraph{Regression tests} \texttt{test\_fee\_fix.py} ensures that a
previous bug (double-counting commission in the itemised view) does not
reappear.

\subsection{Testing Limitations}

There is no end-to-end integration test that checks the full pipeline
(from data ingestion to order execution) in a single run. Backtest mode
partially fills this gap, but a dedicated simulation integration test suite
would improve coverage.
